\subsection{Operations of the Poincaré group and morphisms of coverings}

Let B be a connected and locally path-connected topological space and let b be a point of B.

Let E be a covering of B. Since the space B is assumed to be connected, the fibre \(E_{b}\) is nonempty if E is nonempty (I, p. 74, prop. 4). By assertion a) of theorem 1 of III, p. 305, the covering E is connected and nonempty if and only if the group \(\pi_{1}(B,b)\) acts transitively on \(E_{b}\) .

PROPOSITION 2. --- Let E and E' be coverings of B.

\begin{enumerate}
\def\labelenumi{\alph{enumi})}
\item
  Let \(f\colon\mathrm{E}\to\mathrm{E}^{\prime}\) be a B-morphism. The map \(f_{b}\colon\mathrm{E}_{b}\to\mathrm{E}_{b}^{\prime}\) deduced from \(f\) is compatible with the operations of \(\pi_{1}(\mathrm{B},b)\) on \(\mathrm{E}_{b}\) and \(\mathrm{E}_{b}^{\prime}\) respectively. It is bijective if and only if \(f\) is an isomorphism.
\item
  The map \(f\mapsto f_{b}\) is a bijection from the set \(\mathcal{C}_{\mathrm{B}}(\mathrm{E};\mathrm{E}^{\prime})\) of B-morphisms from E to \(\mathrm{E}^{\prime}\) onto the set \(\mathcal{F}_{\pi_1(\mathrm{B},b)}(\mathrm{E}_b;\mathrm{E}_b^{\prime})\) of \(\pi_1(\mathrm{B},b)\)-morphisms from \(\mathrm{E}_b\) to \(\mathrm{E}_b^{\prime}\) .
\end{enumerate}

If \(f\) is a B-morphism from E to \(\mathrm{E}'\) , the map \(f_b\colon \mathrm{E}_b\to \mathrm{E}_b'\) is a \(\pi_1(\mathrm{B},b)\) -morphism (cf. III, p. 305). Moreover, two B-morphisms from E to \(\mathrm{E}'\) which coincide on the fibre \(\mathrm{E}_b\) are equal (I, p. 80, cor. 3 of prop. 6).

Let \(\varphi\colon\mathrm{E}_{b}\to\mathrm{E}_{b}^{\prime}\) be a morphism of \(\pi_{1}(\mathrm{B},b)\)-sets; let us show that there exists a B-morphism \(f\) from E to \(\mathrm{E}^{\prime}\) such that \(f_{b}=\varphi\) . We may assume the spaces E and \(\mathrm{E}^{\prime}\) are connected and nonempty, so that \(\mathrm{E}_{b}\) and \(\mathrm{E}_{b}^{\prime}\) are homogeneous \(\pi_{1}(\mathrm{B},b)\)-sets. Let \(x\) be a point of \(\mathrm{E}_{b}\) ; its stabilizer G is the image of the map \(\pi_{1}(p,x)\) (III, p. 305, Theorem 1). Since the map \(\varphi\) is a \(\pi_{1}(\mathrm{B},b)\)-morphism, the group G fixes the point \(x^{\prime}=\varphi(x)\) of \(\mathrm{E}_{b}^{\prime}\) , hence is contained in the image of the map \(\pi_{1}(p^{\prime},x^{\prime})\) . By Prop. 1 of III, p. 308, there exists a B-morphism \(f\) from E to \(\mathrm{E}^{\prime}\) such that \(f(x)=x^{\prime}\) . The maps \(f_{b}\) and \(\varphi\) are \(\pi_{1}(\mathrm{B},b)\)-morphisms of the homogeneous \(\pi_{1}(\mathrm{B},b)\)-set \(\mathrm{E}_{b}\) to \(\mathrm{E}_{b}^{\prime}\) which coincide at the point \(x\) ; they are therefore equal.

If \(f\colon \mathrm{E}\to \mathrm{E}'\) is a B-isomorphism, the map \(f_{b}\colon \mathrm{E}_{b}\to \mathrm{E}_{b}^{\prime}\) is bijective. Conversely, suppose the map \(f_{b}\) is bijective and let \(g\colon \mathrm{E}^{\prime}\to \mathrm{E}\) be a B-morphism such that \(g_{b} = (f_{b})^{-1}\) . The B-morphism \(g\circ f\colon \mathrm{E}\to \mathrm{E}\) induces on \(\mathrm{E}_b\) the identity map; it therefore follows from the proposition that \(g \circ f = Id_{E}\) . Likewise, \(f \circ g = Id_{E'}\) , which proves that f is a B-isomorphism.

COROLLARY 1. --- The coverings E and E' are isomorphic if and only if the \(\pi_1(\mathrm{B},b)\) -sets \(\mathrm{E}_b\) and \(\mathrm{E}_b'\) are isomorphic.

COROLLARY 2. --- Let (E,p) and (E',p') be connected coverings of B. Let \(x\) be a point of \(E_{b}\) and \(x'\) a point of \(E'_{b}\) . For there to exist a B-morphism \(g\colon E\to E'\) such that \(g(x)=x'\), it is necessary and sufficient that \(p_{*}(\pi_{1}(E,x))\subset p'_{*}(\pi_{1}(E',x'))\). Such a morphism is then unique and is an isomorphism if and only if the subgroups \(p_{*}(\pi_{1}(E,x))\) and \(p'_{*}(\pi_{1}(E',x'))\) of \(\pi_{1}(B,b)\) are equal.

By III, p. 305, Theorem 1, the map from \(\pi_1(B,b)\) onto \(E_b\) defined by \(\gamma \mapsto x \cdot \gamma\) induces by passage to the quotient an isomorphism of the \(\pi_1(B,b)\)-set \(p_*(\pi_1(E,x)) \backslash \pi_1(B,b)\) onto \(E_b\) . Likewise, there exists a unique isomorphism of \(\pi_1(B,b)\)-sets from \(p'_{*}(\pi_1(E',x')) \backslash \pi_1(B,b)\) onto \(E_b'\) . By Proposition 2, there exists a B-morphism \(g\colon E\to E'\) such that \(g(x)=x'\) if and only if there exists a morphism of \(\pi_1(B,b)\)-sets from \(p_*(\pi_1(E,x)) \backslash \pi_1(B,b)\) onto \(p'_{*}(\pi_1(E',x')) \backslash \pi_1(B,b)\) sending the class \(p_*(\pi_1(E,x))\) to the class \(p'_{*}(\pi_1(E',x'))\) . Such a morphism exists if and only if one has \(p_*(\pi_1(E,x)) \subset p'_{*}(\pi_1(E',x'))\) . It is then unique, since the space E is assumed connected (I, p. 34, Cor. 2 of Prop. 11). It is an isomorphism if and only if the subgroups \(p_*(\pi_1(E,x))\) and \(p'_{*}(\pi_1(E',x'))\) of \(\pi_1(B,b)\) are equal.

COROLLARY 3. --- Let (E,p) be a connected covering of B and let x be a point of the fiber \(E_{b}\) . Let N denote the normalizer of \(p_{*}(\pi_{1}(E,x))\) in \(\pi_{1}(B,b)\) . For every element \(\gamma\) of N, there exists a unique B-automorphism g of E such that \(g(x) = x \cdot \gamma\) , and the map \(\gamma \mapsto g\) defines by passage to the quotient a group isomorphism from \(\mathrm{N}/p_{*}(\pi_{1}(\mathrm{E},x))\) onto \(\mathrm{Aut}_{\mathrm{B}}(\mathrm{E})\) .

Let \(\gamma \in \pi_1(\mathrm{B}, b)\) ; we have \(p_*(\pi_1(\mathrm{E}, x)) = \operatorname{Int}(\gamma)(p_*(\pi_1(\mathrm{E}, x \cdot \gamma)))\) (III, p. 305, Theorem 1, c)). By Corollary 2, for there to exist a B-automorphism \(g\) of E such that \(g(x) = x \cdot \gamma\) , it is necessary and sufficient that \(\gamma\) belongs to N. Such an isomorphism is then unique. Let us denote by \(\alpha: \mathrm{N} \to \mathrm{Aut}_{\mathrm{B}}(\mathrm{E})\) the map \(\gamma \mapsto g\) thus defined.

Let \(\gamma\) and \(\gamma'\) be two elements of N. Set \(g = \alpha(\gamma)\) , \(g' = \alpha(\gamma')\) ; we have \(g(g'(x)) = g(x \cdot \gamma') = g(x) \cdot \gamma' = (x \cdot \gamma) \cdot \gamma' = x \cdot (\gamma\gamma')\) by relations (4), III, p. 305 and (1), p. 304. Consequently, \(\alpha\) is a group homomorphism. For \(\alpha(\gamma) = \mathrm{Id}_{\mathrm{E}}\) , it is necessary and sufficient that one have \(x \cdot \gamma = x\) , by virtue of the uniqueness of \(\alpha(\gamma)\) , that is, \(\gamma \in p_*(\pi_1(E, x))\) (III, p. 305, Theorem 1, b)). Finally, if \(g\) is a B-automorphism of E, there exists a path \(c\) joining \(x\) to \(g(x)\) in E (which is path-connected), and one then has \(g(x) = x \cdot \gamma\) , where \(\gamma\) is the class of the path \(p \circ c\) . This proves that the homomorphism \(\alpha\) is surjective.

The group \(\mathrm{Aut}_{\mathrm{B}}(\mathrm{E})\) acts on the fibre \(\mathrm{E}_b\) and is identified with the group of automorphisms of the homogeneous \(\pi_1(\mathrm{B},b)\) -set (on the right) \(\mathrm{E}_b\) (cf. A, I, p. 56, prop. 5 and 6).

COROLLARY 4. --- Let (E, p) be a connected covering of B and let x be a point of the fibre \(E_{b}\) . For E to be a Galois covering of B, it is necessary and sufficient that \(p_{*}(\pi_{1}(E, x))\) be a distinguished subgroup of \(\pi_{1}(B, b)\) . The group \(\mathrm{Aut}_{\mathrm{B}}(\mathrm{E})\) is then isomorphic to the quotient group \(\pi_{1}(\mathrm{B}, b)/p_{*}(\pi_{1}(\mathrm{E}, x))\) .

If \(p_{*}(\pi_{1}(\mathrm{E},x))\) is a distinguished subgroup of \(\pi_1(\mathrm{B},b)\) , the group \(\mathrm{Aut}_{\mathrm{B}}(\mathrm{E})\) acts transitively on the fibre \(\mathrm{E}_b\) by Corollary 3, so that E is a Galois covering of B (I, p. 102, th. 2). The converse is assertion d) of Theorem 1 (III, p. 305). The last assertion follows from Corollary 3.

